\begin{frame}[label=excel]{5g · Exportar las consultas 5e y 5f con Spark}
  \small
  \pregunta{¿Cómo entregamos los resultados de las seis consultas de vistas temporales en Excel?}
  \flujo{\texttt{spark.sql}\\Consultas 5e y 5f}{DataFrame de resultado\\\texttt{coalesce(1)}}{\texttt{write.format}\\Proveedor spark-excel}{6 libros \texttt{.xlsx}\\Una hoja por consulta}
  \begin{columns}[T,onlytextwidth]
    \column{0.52\textwidth}
    \textbf{Escritura real con Spark}
    \codigo{sql/05g_excel}{102}{110}
    \column{0.45\textwidth}
    \textbf{Salida}\\[5pt]
    \footnotesize
    \begin{tabular}{@{}llr@{}}
      \toprule Consulta & Hoja & Filas\\\midrule
      5e.1 & \texttt{5e1\_variacion} & 60\\
      5e.2 & \texttt{5e2\_perfil} & 18\\
      5e.3 & \texttt{5e3\_region} & 9\\
      5f.1 & \texttt{5f1\_distritos} & 10\\
      5f.2 & \texttt{5f2\_nacionalidad} & 12\\
      5f.3 & \texttt{5f3\_meses} & 12\\\bottomrule
    \end{tabular}
    \par\vspace{5pt}Destino: carpeta \texttt{excel/}.\\
    Proveedor: \texttt{com.crealytics.spark.excel}.
  \end{columns}
  \respuesta{El registro de ejecución documenta seis libros escritos desde DataFrames.
    En 5e.2, el export conserva totales y promedio, pero no incluye la columna de porcentaje nacional de la consulta individual.}
  \origen{sql}{05g_excel}
  \note{40 s. Explicar proveedor, cabecera, hoja y archivo. El código recalcula los seis resultados; no transforma archivos de texto en Excel.}
\end{frame}
